\section{The burden of continual certification}
\label{sec:maintenance}

\subsection{One uncertain arrival can obstruct many reports}

The first lower bound retains one reward environment throughout the time axis.
It therefore measures a cumulative burden without adding the worst cases of
incompatible experiments. Fix
$0<\varepsilon\le1/16$, $0<\delta\le1/16$, and $0<\eta<\varepsilon/2$.
An initial reference has Bernoulli mean $1/4+\varepsilon/2$. Candidates arrive
on any prescribed schedule. In environment zero each has mean $p=1/4$;
environment $j$ changes only candidate $j$ to mean $p+2\varepsilon$.
All candidates are good in environment zero, with boundary distance
$\varepsilon/2$. After candidate $j$ arrives in environment $j$, every other
candidate is bad.

Let $M_t$ count arrived candidates and $a_t=M_t-M_{t-1}$. Define
\begin{align}
 q_-&=\left\lfloor\frac{\kl(1/4,\delta)}
                              {\kl(1/4,1/4+2\varepsilon)}\right\rfloor,
 &Q_t&=(Q_{t-1}+q_-a_t-1)_+,\quad Q_0=0,\notag\\
 \mathcal B_T&=\sum_{t\le T}(M_t-1)_+\ind\{Q_t>0\}.
 \label{eq:broadcast_queue}
\end{align}
The queue is an information-capacity certificate: each arriving candidate
requires $q_-$ units, while only one observation is possible per round. It is
not the queue of any particular policy. In the stated range $q_-\ge2$.

\begin{theorem}[Broadcast uncertainty]
\label{thm:broadcast}
For every policy reliable on environment zero and all its one-candidate
alternatives,
\[
 \E_0\mathcal A_T(\eta)\ge(3/4-\delta)\mathcal B_T.
\]
If $\mathcal B_T>0$, then
\[
 \Prob_0\{\mathcal A_T(\eta)\ge\mathcal B_T/4\}\ge1/4.
\]
These conclusions still hold when the policy knows the full arrival schedule,
the horizon, and the reference distribution.
\end{theorem}

At a busy time, some deterministic candidate label has expected sample count
below $q_-$. Changing that candidate alone has too little history KL to permit
positive certification of another candidate with probability at least $1/4$.
Reliability also bounds incorrect negative reports. The remaining probability
must be unresolved. The same reference environment supports this argument at
every time; the difficult label is selected from expected counts rather than
from a favorable realized history. The supplementary proof includes the
bounded-variable step needed for the probability statement.

If $m$ candidates have arrived by time $s$, the queue obeys the deterministic
inequality
\begin{equation}
 \sum_{t\le s+qm-1}(M_t-1)_+\ind\{Q_t(q)>0\}
                  \ge\frac{q-1}{2}m(m-1),\qquad q\ge2.
 \label{eq:queue_area}
\end{equation}
The factor $m^2$ reflects the number of affected reports as well as the work
needed to exclude high-quality alternatives. The subtraction of one unit in
$q-1$ accounts for the last observation of each busy period: reports are issued
after service.

\begin{corollary}[Arrival-growth boundary]
\label{cor:growth}
Fix the preceding task parameters and one infinite arrival schedule. If a
policy is reliable for every reward environment and has a deterministic
envelope $a_T=o(T)$ satisfying \eqref{eq:joint_protocol} in every environment,
then $K_t=o(\sqrt t)$. Conversely, $K_T^2\log(eK_T)=o(T)$ suffices for one
horizon-free policy to have both $\E R_T=o(T)$ and such a sublinear area
envelope. In particular, for $K_t\asymp t^\alpha$, $0<\alpha<1$, simultaneous
attainment is possible if and only if $\alpha<1/2$.
\end{corollary}

The sufficiency follows from the next theorem. The exponent statement leaves
a logarithmic gap for general critical schedules. It does not identify
$K_t=o(\sqrt t)$ as a sufficient condition on every schedule.

\subsection{Worst-schedule maintenance area}

Let $\mathfrak A_{T,K}$ be the smallest $a$ attainable by a full-model reliable
policy with
\[
 \sup_{\tau:K_T\le K}\sup_\mu
 \Prob_\mu\{E_{\rm lab}^c\ \text{or}\ \mathcal A_T(\eta)>a\}\le\delta.
\]
Policies may know $T,K$ and the schedule for the lower bound. The upper bound
uses neither $T$ nor $K$ as input.

\begin{theorem}[Worst-schedule area and joint attainability]
\label{thm:worst_area}
For $K\ge2$, $T\ge1$, $0<\varepsilon\le1/8$,
$0<\eta\le\varepsilon/32$, and $0<\delta\le1/16$,
\begin{equation}
 2^{-24}\min\left\{KT,\frac{K^2}{\eta^2}\log\frac K\delta\right\}
 \le\mathfrak A_{T,K}\le
 256\min\left\{KT,\frac{K^2}{\eta^2}\log\frac K\delta\right\}.
 \label{eq:worst_area}
\end{equation}
The CP policy with $\rho=\eta/4$ attains the upper order and also satisfies
\begin{equation}
 \E R_T\le\min\left\{T,K+4\sqrt{3KT\log(2T)}+1+
                         99K\eta^{-2}\log(K/\delta)\right\}.
 \label{eq:joint_general}
\end{equation}
\end{theorem}

For the upper bound, completed checkpoints imply endpoint errors at most
$\eta/4$, sufficient to classify all arms outside the boundary band. The
unfinished-work queue is dominated by adding $q_i$ at each arrival and serving
at even rounds. On the simultaneous coverage event,
\begin{equation}
 \mathcal A_T(\eta)\le\min\left\{K_TT,
                         2K_T\sum_{i\le K_T}q_i\right\}.
 \label{eq:queue_upper}
\end{equation}
Additional UCB observations reduce the work; they are not charged again as
new checkpoint samples. The lower bound uses $K-1$ candidates with independent
means $1/2\pm2\eta$ and a reference of mean $1/2+\varepsilon$, all available
at the first pull. Simultaneous partial recovery and a horizon truncation give
the probability statement required by the definition of $\mathfrak A_{T,K}$.
The statistical $\log K$ factor has fixed-arm multiple-testing precedents;
the lower construction itself requires no repeated arrivals.

For fixed task parameters, \eqref{eq:worst_area} implies a precise uniform
terminal-count statement: over all schedules with $K_T\le k(T)$, sublinear
expected regret and sublinear service envelopes are jointly attainable if and
only if $k(T)^2\log k(T)=o(T)$. Here the worst schedule can change with $T$.
On a single fixed schedule that conclusion cannot be inferred from total
growth alone. For example, batches of size
\[
 m_n=\left\lfloor\frac{\sqrt{s_n}}{(\log s_n)^{1/4}}\right\rfloor,
 \qquad s_n=2^{2^n},
\]
give $K_t=o(\sqrt t)$ but preclude any $o(T)$ area envelope uniform over
reward environments. Near $s_n$, simultaneous classification of that burst
requires accumulated unresolved area of order $m_n^2\log m_n$, which exceeds
$s_n$ by a growing factor. The supplementary argument selects a single
infinite schedule; the difficult reward environment is allowed to vary with
the endpoint, as required by the stated uniform guarantee.
